\chapter{Dokumentasi Praktikum}
\section{Hasil Praktikum}
\subsection{Launch Unreal Engine}

\begin{figure}[H]
    \centering
    \includegraphics[width=0.6\textwidth]{1.png}
    \caption{Proses Instalasi Unreal Engine di Epic Games Launcher}
    \label{fig:1}
\end{figure}
Keterangan Gambar :
\begin{itemize}
    \item Buka aplikasi Epic Games Launcher dan navigasikan ke tab \textbf{Library}.
    \item Lakukan instalasi Unreal Engine versi 5.8.2 dan tunggu hingga proses pengunduhan selesai.
\end{itemize}

\begin{figure}[H]
    \centering
    \includegraphics[width=0.6\textwidth]{2.png}
    \caption{Pembuatan Proyek Baru di Unreal Engine}
    \label{fig:2}
\end{figure}
Keterangan Gambar :
\begin{itemize}
    \item Buka Unreal Engine yang telah terinstal untuk memunculkan jendela \textbf{Unreal Project Browser}.
    \item Pilih tab kategori \textbf{GAMES}, lalu klik template \textbf{Intro To Unreal}.
    \item Tentukan lokasi penyimpanan dan masukkan nama proyek yang diinginkan, kemudian klik tombol \textbf{Create}.
\end{itemize}

\begin{figure}[H]
    \centering
    \includegraphics[width=0.6\textwidth]{3.png}
    \caption{Splash Screen Memuat Unreal Editor}
    \label{fig:3}
\end{figure}
Keterangan Gambar :
\begin{itemize}
    \item Tunggu proses inisialisasi pada splash screen saat memuat proyek ke dalam Unreal Editor.
    \item Pastikan proses \textit{Compiling blueprints} berjalan hingga selesai (100\%) dan editor terbuka sepenuhnya.
\end{itemize}

\begin{figure}[H]
    \centering
    \includegraphics[width=0.6\textwidth]{4.png}
    \caption{Antarmuka Utama Unreal Editor}
    \label{fig:4}
\end{figure}
Penjelasan Antarmuka Utama Unreal :
\begin{itemize}
    \item Tinjau bagian \textbf{Documentation Page} di sisi kiri; jendela ini berisi panduan dan bisa ditutup bila sedang tidak diperlukan.
    \item Gunakan area \textbf{Viewport} di bagian tengah sebagai layar utama untuk melihat dan bernavigasi di dalam level 3D.
    \item Perhatikan panel \textbf{Outliner} di sisi kanan atas yang berfungsi untuk menampilkan daftar semua objek yang ada di dalam level.
    \item Gunakan panel \textbf{Details} di sisi kanan bawah untuk melihat dan mengubah \textit{properties} dari objek dalam level yang sedang dipilih/klik.
\end{itemize}

\subsection{Eksplorasi Antarmuka dan Fungsionalitas Dasar (Materi)}

\subsubsection{Fungsi Top Bar}
\begin{figure}[H]
    \centering
    \includegraphics[width=0.6\textwidth]{23.png}
    \caption{Tampilan Fungsi Top Bar}
    \label{fig:23}
\end{figure}
Penjelasan :
\begin{itemize}
    \item \textbf{Deskripsi:} Menampilkan bilah menu utama (\textit{top bar}) pada antarmuka editor Unreal Engine.
    \item \textbf{Keterangan:} Bagian ini bersifat informatif, menyajikan pintasan menu esensial seperti tombol \textit{Save}, \textit{Selection Mode}, tombol \textit{Play} untuk menguji level, hingga akses menuju \textit{Blueprints Menu} dan \textit{Cinematics/Sequencer}.
\end{itemize}

\subsubsection{Pengaturan Grid}
\begin{figure}[H]
    \centering
    \includegraphics[width=0.6\textwidth]{24.png}
    \caption{Pengaturan Snap ke Grid}
    \label{fig:24}
\end{figure}
Penjelasan :
\begin{itemize}
    \item \textbf{Deskripsi:} Fitur penyesuaian skala \textit{grid} atau batas spasial pada \textit{viewport}.
    \item \textbf{Keterangan:} Informasi ini menunjukkan letak ikon pengatur jarak \textit{grid} yang memastikan objek akan otomatis menempel (\textit{snap}) secara presisi ketika dipindahkan.
\end{itemize}

\subsubsection{Transformasi Objek}
\begin{figure}[H]
    \centering
    \includegraphics[width=0.6\textwidth]{5.png}
    \caption{Transformasi Objek - Move}
    \label{fig:5}
\end{figure}

\begin{figure}[H]
    \centering
    \includegraphics[width=0.6\textwidth]{6.png}
    \caption{Transformasi Objek - Rotate}
    \label{fig:6}
\end{figure}

\begin{figure}[H]
    \centering
    \includegraphics[width=0.6\textwidth]{7.png}
    \caption{Transformasi Objek - Scale}
    \label{fig:7}
\end{figure}
Penjelasan :
\begin{itemize}
    \item \textbf{Deskripsi:} Mengilustrasikan instrumen transformasi (\textit{gizmo}) dasar untuk mengedit properti objek.
    \item \textbf{Keterangan:} Memperlihatkan secara terpisah bentuk sumbu panah untuk alat \textit{Move} (menggeser koordinat), lingkaran untuk \textit{Rotate} (memutar sudut objek), dan kubus ujung untuk \textit{Scale} (mengubah skala ukuran). 
\end{itemize}

\subsubsection{Content Drawer (Pencarian Aset)}
\begin{figure}[H]
    \centering
    \includegraphics[width=0.6\textwidth]{8.png}
    \caption{Pencarian Objek Block Melalui Content Drawer}
    \label{fig:8}
\end{figure}
Penjelasan :
\begin{itemize}
    \item \textbf{Deskripsi:} Penggunaan panel \textit{Content Drawer} sebagai pustaka pengambilan aset.
    \item \textbf{Keterangan:} Pada tahap materi ini, terlihat bahwa kondisi area lingkungan (\textit{level}) belum memiliki jalan penghubung (terputus), sehingga pemain terisolasi. Oleh karena itu, kita perlu mencari dan mengambil aset \textit{block} baru dari \textit{Content Drawer}.
\end{itemize}

\subsubsection{Details Panel dan Collision (Masalah Tembus Ruang)}
\begin{figure}[H]
    \centering
    \includegraphics[width=0.6\textwidth]{9.png}
    \caption{Tampilan Awal Tanjakan (Ramp) Tanpa Collision}
    \label{fig:9}
\end{figure}

\begin{figure}[H]
    \centering
    \includegraphics[width=0.6\textwidth]{10.png}
    \caption{Proses Pengaturan Collision pada Details Panel}
    \label{fig:10}
\end{figure}
Penjelasan :
\begin{itemize}
    \item \textbf{Deskripsi:} Pengaturan awal pada sebuah objek tanjakan (\textit{ramp}) sebelum ditambahkan fungsi batas benturan fisik (\textit{Collision}).
    \item \textbf{Keterangan:} Karena objek \textit{ramp} masih memiliki pengaturan \textit{default} (tidak ada \textit{collision}), objek ini tidak memiliki pijakan padat. Jika pemain menabrak atau mencoba menaikinya, karakter akan otomatis *nembus* jatuh ke dalam \textit{ramp} tersebut.
\end{itemize}

\subsubsection{Outliner (Rintangan Jalan)}
\begin{figure}[H]
    \centering
    \includegraphics[width=0.6\textwidth]{11.png}
    \caption{Pemilihan Objek Pembatas yang Menghalangi Lewat Outliner}
    \label{fig:11}
\end{figure}
Penjelasan :
\begin{itemize}
    \item \textbf{Deskripsi:} Tampilan objek pembatas fisik (\textit{barrier}) di panel \textit{Outliner}.
    \item \textbf{Keterangan:} Pada kondisi awal ini, karakter pemain terblokir dan tidak bisa lewat untuk melanjutkan perjalanan karena terdapat rintangan yang melintang di jalan. Objek penghalang ini sedang diidentifikasi dan disorot melalui panel \textit{Outliner}.
\end{itemize}

\subsubsection{Materials (Material Bawaan)}
\begin{figure}[H]
    \centering
    \includegraphics[width=0.6\textwidth]{12.png}
    \caption{Kondisi Objek Masih dengan Material Default}
    \label{fig:12}
\end{figure}
Penjelasan :
\begin{itemize}
    \item \textbf{Deskripsi:} Tampilan dasar suatu objek sebelum dihias dengan tekstur permukaan.
    \item \textbf{Keterangan:} Objek tersebut masih menggunakan material bawaan (\textit{default material}) dari mesin Unreal Engine, sehingga hanya tampak datar, abu-abu polos, dan tidak memancarkan kesan estetika apa pun.
\end{itemize}

\subsubsection{Particles Effect (Intensitas Rendah)}
\begin{figure}[H]
    \centering
    \includegraphics[width=0.6\textwidth]{13.png}
    \caption{Kondisi Awal Pancaran Partikel (Counts Masih Sedikit)}
    \label{fig:13}
\end{figure}
Penjelasan :
\begin{itemize}
    \item \textbf{Deskripsi:} Konfigurasi parameter pancaran awal pada \textit{emitter} sistem partikel.
    \item \textbf{Keterangan:} Pada sesi materi, terlihat bahwa parameter \textit{counts} (jumlah partikel yang dilepaskan dalam satu siklus) masih sangat sedikit, sehingga efek animasi visual yang dihasilkan terlihat kosong atau tidak pekat.
\end{itemize}

\subsubsection{MetaSounds (Random Seed Default)}
\begin{figure}[H]
    \centering
    \includegraphics[width=0.6\textwidth]{14.png}
    \caption{Pengaturan Awal MetaSounds dengan Random Seed Bawaan}
    \label{fig:14}
\end{figure}
Penjelasan :
\begin{itemize}
    \item \textbf{Deskripsi:} Kondisi nilai parameter acak generator suara pada sistem MetaSounds.
    \item \textbf{Keterangan:} Node pembentuk suara prosedural masih menggunakan nilai \textit{random seed} yang murni bawaan (\textit{default}), sebelum dilakukannya modifikasi atau penyetelan ulang nilai variabel.
\end{itemize}

\subsubsection{Sistem Fisika (Manipulasi Massa Objek)}
\begin{figure}[H]
    \centering
    \includegraphics[width=0.6\textwidth]{15.png}
    \caption{Pengaturan Massa Objek Menjadi Lebih Ringan}
    \label{fig:15}
\end{figure}
Penjelasan :
\begin{itemize}
    \item \textbf{Deskripsi:} Panel modifikasi besaran bobot (\textit{mass}) fisika objek yang ada di properti Detail.
    \item \textbf{Keterangan:} Tampilan awal ini mendemonstrasikan proses manipulasi hukum alam di mana nilai angka massa dari sebuah beban dikonfigurasi dan direkayasa agar menjadi jauh lebih ringan dibanding bentuk atau ukuran aslinya.
\end{itemize}

\subsubsection{Pencahayaan (Lighting)}
\begin{figure}[H]
    \centering
    \includegraphics[width=0.6\textwidth]{16.png}
    \caption{Kondisi Pencahayaan Level Bawaan (Default)}
    \label{fig:16}
\end{figure}
Penjelasan :
\begin{itemize}
    \item \textbf{Deskripsi:} Representasi pencahayaan ruangan tingkat standar.
    \item \textbf{Keterangan:} Lingkungan masih terlihat datar dengan distribusi cahaya yang seadanya, membuat dimensi ruang serta kedalaman rona bayangan menjadi kurang terasa (\textit{flat}).
\end{itemize}

\subsubsection{Animasi (Materi Informatif)}
\begin{figure}[H]
    \centering
    \includegraphics[width=0.6\textwidth]{17.png}
    \caption{Antarmuka Jendela Peninjauan Animasi}
    \label{fig:17}
\end{figure}
Penjelasan :
\begin{itemize}
    \item \textbf{Deskripsi:} Jendela manajemen aset animasi untuk karakter berkerangka (\textit{skeletal mesh}).
    \item \textbf{Keterangan:} Gambar memamerkan area di mana kita bisa melihat pratinjau pergerakan karakter (seperti berjalan atau diam) yang bisa dipasangkan untuk membangun interaksi gerak yang realistis.
\end{itemize}

\subsubsection{Sequencer (Persiapan Jalur Kamera)}
\begin{figure}[H]
    \centering
    \includegraphics[width=0.6\textwidth]{18.png}
    \caption{Materi Pengenalan Timeline Sequencer}
    \label{fig:18}
\end{figure}
Penjelasan :
\begin{itemize}
    \item \textbf{Deskripsi:} Tahap inisialisasi awal (*materi*) antarmuka Sequencer sebelum dilakukan aksi perekaman.
    \item \textbf{Keterangan:} Jendela \textit{timeline} sedang dikosongkan dan dipersiapkan agar parameter adegan dan pergerakan lintasan kamera (*camera track*) bisa disematkan ke dalamnya.
\end{itemize}

\subsubsection{Sistem Blueprint (Pintu Konfeti dan Error Lampu)}
\begin{figure}[H]
    \centering
    \includegraphics[width=0.6\textwidth]{19.png}
    \caption{Interaksi Ketika Button Ditekan (Lampu Atas Masih Merah)}
    \label{fig:19}
\end{figure}

\begin{figure}[H]
    \centering
    \includegraphics[width=0.6\textwidth]{20.png}
    \caption{Penyusunan Logika untuk Memperbaiki Bug Blueprint}
    \label{fig:20}
\end{figure}
Penjelasan :
\begin{itemize}
    \item \textbf{Deskripsi:} Eksekusi awal Blueprint pada skenario pembukaan pintu masuk.
    \item \textbf{Keterangan:} Saat tombol ditekan, pintu mekanis terbuka disambut letupan konfeti (*confetti*) dan lampu indikator kanan-kiri sukses berubah hijau. Sayangnya, terdapat \textit{bug}: lampu bagian atas masih dibaca salah oleh Blueprint sehingga tetap menyala merah. Gambar selanjutnya menunjukkan upaya perbaikan (*setting*) koneksi *node* Blueprint.
\end{itemize}

\subsubsection{Panduan Lanjutan / Next Steps (Materi Informatif)}
\begin{figure}[H]
    \centering
    \includegraphics[width=0.6\textwidth]{37.png}
    \caption{Halaman Dokumentasi Panduan Lanjutan (Next Steps)}
    \label{fig:37}
\end{figure}
Penjelasan :
\begin{itemize}
    \item \textbf{Deskripsi:} Dokumentasi portal rujukan pengembang perangkat lunak dari Epic Games.
    \item \textbf{Keterangan:} Rangkuman arahan fitur-fitur apa lagi yang dapat diunduh dan dipelajari untuk memperdalam penguasaan mesin Unreal Engine setelah pemahaman materi dasar selesai.
\end{itemize}

\subsubsection{Templates Proyek (Materi Informatif)}
\begin{figure}[H]
    \centering
    \includegraphics[width=0.6\textwidth]{21.png}
    \caption{Berbagai Pilihan Base Template Permainan}
    \label{fig:21}
\end{figure}
Penjelasan :
\begin{itemize}
    \item \textbf{Deskripsi:} Menu ketika pertama kali membuat proyek instalasi Unreal Engine.
    \item \textbf{Keterangan:} Menyajikan pilihan kerangka fungsional (\textit{First-Person}, \textit{Third-Person}, dll.) untuk mempercepat fase pembuatan game tanpa perlu merakit elemen mekanik gerak dari nol.
\end{itemize}

\subsubsection{Samples Proyek (Materi Informatif)}
\begin{figure}[H]
    \centering
    \includegraphics[width=0.6\textwidth]{22.png}
    \caption{Daftar Sampel Proyek yang Tersedia}
    \label{fig:22}
\end{figure}
Penjelasan :
\begin{itemize}
    \item \textbf{Deskripsi:} Menu pameran (*showcase*) proyek sampel jadi berskala besar.
    \item \textbf{Keterangan:} Menyediakan aset lingkungan kualitas tinggi yang disediakan khusus oleh Epic Games sebagai studi kasus terbaik dalam praktik pengembangan game.
\end{itemize}
\clearpage
\section{Hasil Akhir}

Bagian ini memaparkan hasil fungsional maupun modifikasi perbaikan paska eksekusi dari skenario kendala materi-materi di atas:

\subsection{Hasil Transformasi Objek Terhadap Lingkungan}
\begin{figure}[H]
    \centering
    \includegraphics[width=0.6\textwidth]{25.png}
    \caption{Hasil Peletakan dan Skala Objek pada Area}
    \label{fig:25}
\end{figure}
Penjelasan :
\begin{itemize}
    \item \textbf{Deskripsi:} Hasil perpaduan akhir bentuk tata letak objek.
    \item \textbf{Keterangan:} Membuktikan bahwa pemakaian fungsi Move, Rotate, dan Scale membuahkan hasil penyesuaian yang proporsional pada objek sehingga menyatu dengan levelnya secara utuh.
\end{itemize}

\subsection{Jembatan Penyeberangan Lewat Aset Content Drawer}
\begin{figure}[H]
    \centering
    \includegraphics[width=0.6\textwidth]{26.png}
    \caption{Hasil Penambahan Objek Block Sebagai Jalan Keluar}
    \label{fig:26}
\end{figure}
Penjelasan :
\begin{itemize}
    \item \textbf{Deskripsi:} Kondisi area pasca penyelesaian jalan yang sebelumnya terputus.
    \item \textbf{Keterangan:} Menunjukkan permasalahan isolasi jalur telah terpecahkan. Menggunakan metode \textit{drag and drop}, objek baru jenis kotak (\textit{block}) yang diambil dari \textit{Content Drawer} sukses dipasang sebagai jembatan yang memungkinkan karakter akhirnya bisa menyeberang.
\end{itemize}

\subsection{Implementasi Fisik Collision pada Tanjakan (Ramp)}
\begin{figure}[H]
    \centering
    \includegraphics[width=0.6\textwidth]{27.png}
    \caption{Karakter Sukses Berpijak pada Ramp Tanpa Menembus}
    \label{fig:27}
\end{figure}
Penjelasan :
\begin{itemize}
    \item \textbf{Deskripsi:} Resolusi konflik kontak fisik permukaan aset.
    \item \textbf{Keterangan:} Jika sebelumnya pemain celaka akibat "*nembus*" langsung ke dalam model grafis \textit{ramp}, berkat implementasi pengaturan geometri \textit{Collision}, tanjakan tersebut kini berubah padat. Pemain dapat menabrak dan melangkah mulus menaiki permukaan tersebut hingga ke puncak.
\end{itemize}

\subsection{Bebas Melintas Usai Pembersihan via Outliner (Clean Blocking Way)}
\begin{figure}[H]
    \centering
    \includegraphics[width=0.6\textwidth]{28.png}
    \caption{Jalanan Area Level Terbuka Paska Penghapusan Rintangan}
    \label{fig:28}
\end{figure}
Penjelasan :
\begin{itemize}
    \item \textbf{Deskripsi:} Kondisi jalur lintasan karakter sesudah objek palang disingkirkan.
    \item \textbf{Keterangan:} Objek *barrier* penghalang yang sebelumnya menghadang tepat di tengah lintasan dan dipilih lewat Outliner, kini telah berhasil dihapus secara bersih (\textit{delete}). Pemain kini leluasa melewati ruang tersebut tanpa terblokir lagi.
\end{itemize}

\subsection{Eksperimen Berbagai Tipe Materials pada Objek}
\begin{figure}[H]
    \centering
    \includegraphics[width=0.6\textwidth]{29.png}
    \caption{Hasil Percobaan Menggunakan Berbagai Material Berbeda}
    \label{fig:29}
\end{figure}
Penjelasan :
\begin{itemize}
    \item \textbf{Deskripsi:} Visual aset objek pasca pengaplikasian lapisan eksperimental.
    \item \textbf{Keterangan:} Transformasi penampilannya sungguh signifikan. Objek melepaskan belenggu material \textit{default}-nya yang kusam dan kini tampak begitu realistis setelah dieksekusi menggunakan percobaan dengan berbagai tipe material pelapis (\textit{surface textures}).
\end{itemize}

\subsection{Kepadatan Penuh dari Particle Counts Sebesar 455}
\begin{figure}[H]
    \centering
    \includegraphics[width=0.6\textwidth]{30.png}
    \caption{Pancaran Partikel Meriah (Total Particle Counts: 455)}
    \label{fig:30}
\end{figure}
Penjelasan :
\begin{itemize}
    \item \textbf{Deskripsi:} Visualisasi luapan efek pancaran \textit{emitter} partikel yang dijalankan (*Play*).
    \item \textbf{Keterangan:} Berbanding terbalik dengan intensitas rendah di awal materi, gambar pasca modifikasi ini mengilustrasikan partikel dalam jumlah (*counts*) yang sangat agresif, yakni di angka \textbf{455}. Hal tersebut menciptakan lautan efek yang meriah dan masif secara visual.
\end{itemize}

\subsection{Keberagaman Audio Berkat Perubahan Random Seed (MetaSounds)}
\begin{figure}[H]
    \centering
    \includegraphics[width=0.6\textwidth]{31.png}
    \caption{Modifikasi Parameter Random Seed pada Graf}
    \label{fig:31}
\end{figure}

\begin{figure}[H]
    \centering
    \includegraphics[width=0.6\textwidth]{32.png}
    \caption{Hasil Keluaran Suara Prosedural yang Lebih Bervariasi}
    \label{fig:32}
\end{figure}
Penjelasan :
\begin{itemize}
    \item \textbf{Deskripsi:} Hasil perombakan generator modulasi bilangan acak pada node MetaSounds.
    \item \textbf{Keterangan:} Dengan ditanggalkannya nilai bawaan (\textit{default seed}), node generator sistem suara prosedural mulai bekerja mengacak keluaran nada. Hasil perubahannya membuktikan bahwa variasi suara yang diproduksi saat \textit{game} dijalankan berubah-ubah tidak monoton.
\end{itemize}

\subsection{Respons Fisika dari Tabrakan pada Massa Ringan}
\begin{figure}[H]
    \centering
    \includegraphics[width=0.6\textwidth]{33.png}
    \caption{Posisi Objek Tergeser Akibat Ditabrak oleh Pemain}
    \label{fig:33}
\end{figure}
Penjelasan :
\begin{itemize}
    \item \textbf{Deskripsi:} Perilaku fisika riil pada benda ringan ketika dipicu oleh persinggungan.
    \item \textbf{Keterangan:} Sesuai dengan pengurangan angka beban (\textit{mass}) yang dilakukan pada fase materi, hasil jadinya tervalidasi sepenuhnya saat objek ditabrak oleh fisik karakter pemain. Posisi benda tersebut bereaksi dengan cepat (terdorong bergeser/terpental), mengamini bobot ringannya di dalam sistem kalkulasi *physics*.
\end{itemize}

\subsection{Dramatisasi Lingkungan Melalui Penyesuaian Pencahayaan (Lighting)}
\begin{figure}[H]
    \centering
    \includegraphics[width=0.6\textwidth]{34.png}
    \caption{Ruangan Sinematik Paska Penempatan Titik Pencahayaan}
    \label{fig:34}
\end{figure}
Penjelasan :
\begin{itemize}
    \item \textbf{Deskripsi:} Visualisasi atmosfer ruang setelah penambahan titik cahaya buatan.
    \item \textbf{Keterangan:} Tidak lagi datar maupun statis, lingkungan bermain dibalut nuansa terang benderang dengan jatuhan rona siluet bayangan dinamis, menambah kental dimensi kedalaman visual layaknya adegan film.
\end{itemize}

\subsection{Pergerakan Kamera Sinematik (Menjalankan Sequencer)}
\begin{figure}[H]
    \centering
    \includegraphics[width=0.6\textwidth]{35.png}
    \caption{Tampilan Lensa Kamera yang Sedang Bergerak Menyusuri Sequencer}
    \label{fig:35}
\end{figure}
Penjelasan :
\begin{itemize}
    \item \textbf{Deskripsi:} Uji coba eksekusi pemutaran jalur \textit{cutscene} di \textit{Sequencer}.
    \item \textbf{Keterangan:} Alih-alih kosong dan diam, status alat Sequencer ini memperagakan contoh rekaman adegan (\textit{in-action}). Sorotan titik fokus kamera beranjak melintasi level, bergerak secara konstan menyesuaikan dengan \textit{timeline keyframe} yang sudah dirancang sebelumnya.
\end{itemize}

\subsection{Penyelesaian Perbaikan Bug Indikator Lampu (Blueprint)}
\begin{figure}[H]
    \centering
    \includegraphics[width=0.6\textwidth]{36.png}
    \caption{Tombol Ditekan, Seluruh Lampu Pintu Normal (Hijau)}
    \label{fig:36}
\end{figure}
Penjelasan :
\begin{itemize}
    \item \textbf{Deskripsi:} Keadaan fungsional paripurna saat eksekusi sambungan node \textit{Blueprint} permainan.
    \item \textbf{Keterangan:} \textit{Bug} malfungsi kelistrikan indikator telah teratasi. Sesudah parameter logika algoritma pada graf disetel ulang (\textit{setting fix}), skenario interaksi merespons dengan benar: begitu *button* terinjak, secara bersamaan taburan konfeti jatuh disusul perubahan ketiga bohlam lampu (termasuk lampu bagian atas) menjadi hijau segar seutuhnya.
\end{itemize}

\chapter{Daftar Pustaka}
\begin{enumerate}
    \item \label{gregory} J. Gregory. \textit{Game Engine Architecture}, 3rd ed. Boca Raton, FL: CRC Press, 2018.
    
    \item \label{romero} M. Romero. \textit{Blueprints Visual Scripting for Unreal Engine 5}, 3rd ed. Birmingham, UK: Packt Publishing, 2022.
    
    \item \label{korkmaz} O. Korkmaz. \textit{Mastering Unreal Engine 4.X: A professional-level guide to building stunning, high-quality games}. Birmingham, UK: Packt Publishing, 2021.
    
    \item \label{akenine} T. Akenine-Möller, E. Haines, dan N. Hoffman. \textit{Real-Time Rendering}, 4th ed. Boca Raton, FL: A K Peters/CRC Press, 2018.
    
    \item \label{dickinson} M. Dickinson. \textit{Unreal Engine 4 Virtual Reality Projects: Build immersive, real-world VR applications}. Birmingham, UK: Packt Publishing, 2019.
    
    \item \label{epic} Epic Games. \textit{Unreal Engine 5 Documentation}. [Online]. Tersedia: \url{https://docs.unrealengine.com/}.
    
    \item \label{sewell} B. Sewell. \textit{Unreal Engine 4 Game Development Quick Start Guide}. Birmingham, UK: Packt Publishing, 2019.
    
    \item \label{dorantes} M. Dorantes. \textit{Unreal Engine 4 for Beginners: Learn to create games with Blueprint}. Independently published, 2020.
\end{enumerate}